% ============================================================================
% CAPÍTULO 4: JUSTIFICACIÓN
% ============================================================================
%
% REQUISITOS SEGÚN LA FACULTAD:
%
% Redacta la justificación de TU INVESTIGACIÓN (no del tema).
% Considera los siguientes criterios:
%
% 1. CONVENIENCIA: ¿Qué tan conveniente es la investigación?; esto es, ¿para qué sirve?
%
% 2. RELEVANCIA SOCIAL: ¿Cuál es su trascendencia para la sociedad?, ¿quiénes se
%    beneficiarán con los resultados de la investigación?, ¿de qué modo? En resumen,
%    ¿qué alcance o proyección social tiene?
%
% 3. IMPLICACIONES PRÁCTICAS: ¿Ayudará a resolver algún problema real?, ¿tiene
%    implicaciones trascendentales para una amplia gama de problemas prácticos?
%
% 4. VALOR TEÓRICO: Con la investigación, ¿se llenará algún vacío de conocimiento?,
%    ¿se podrán generalizar los resultados a principios más amplios?, ¿la información
%    que se obtenga puede servir para revisar, desarrollar o apoyar una teoría?,
%    ¿se podrá conocer en mayor medida el comportamiento de una o de diversas
%    variables o la relación entre ellas?, ¿se ofrece la posibilidad de una exploración
%    fructífera de algún fenómeno o ambiente?, ¿qué se espera saber con los resultados
%    que no se sabía antes?, ¿se pueden sugerir ideas, recomendaciones o hipótesis
%    para futuros estudios?
%
% 5. UTILIDAD METODOLÓGICA: ¿La investigación puede ayudar a crear un nuevo
%    instrumento para recolectar o analizar datos?, ¿contribuye a la definición de
%    un concepto, variable o relación entre variables?, ¿pueden lograrse con ella
%    mejoras en la forma de experimentar con una o más variables?, ¿sugiere cómo
%    estudiar más adecuadamente una población?
% ============================================================================

\chapter{Justificación}
\label{cap:justificacion}

% INSTRUCCIONES GENERALES:
% - Justifica tu INVESTIGACIÓN, no el tema en general
% - Responde ESPECÍFICAMENTE a cada uno de los 5 criterios
% - Sé concreto: menciona quiénes, cómo, cuándo, dónde
% - Conecta con tu problema, pregunta y objetivos
% - Extensión típica: 5-8 páginas

[Párrafo introductorio explicando por qué es importante realizar ESTA investigación específica]

La realización de esta investigación se justifica por las siguientes razones:

% ----------------------------------------------------------------------------
% CRITERIO 1: CONVENIENCIA
% ----------------------------------------------------------------------------
\section{Conveniencia}
\label{sec:conveniencia}

% Pregunta clave: ¿Para qué sirve esta investigación?

Esta investigación es conveniente porque \textit{[explica para qué sirve específicamente]}.

\textbf{Utilidad específica:}

\begin{itemize}
    \item Permitirá \textit{[utilidad 1]}
    \item Proporcionará \textit{[utilidad 2]}
    \item Facilitará \textit{[utilidad 3]}
\end{itemize}

Los resultados de este estudio servirán para \textit{[aplicación concreta]}, lo cual es necesario porque \textit{[razón específica]}.

% ----------------------------------------------------------------------------
% CRITERIO 2: RELEVANCIA SOCIAL
% ----------------------------------------------------------------------------
\section{Relevancia Social}
\label{sec:relevancia_social}

% Preguntas clave:
% - ¿Cuál es su trascendencia para la sociedad?
% - ¿Quiénes se beneficiarán?
% - ¿De qué modo?
% - ¿Qué alcance o proyección social tiene?

\subsection{Trascendencia Social}

Esta investigación tiene relevancia social porque \textit{[explica la trascendencia]}.

\subsection{Beneficiarios Directos}

Los beneficiarios directos de esta investigación son:

\begin{enumerate}
    \item \textbf{[Grupo beneficiario 1]:} Se beneficiarán mediante \textit{[cómo específicamente]}. Se estima que \textit{[número/proporción de personas]} podrían beneficiarse.
    
    \item \textbf{[Grupo beneficiario 2]:} Obtendrán \textit{[beneficio específico]} a través de \textit{[mecanismo]}.
    
    \item \textbf{[Grupo beneficiario 3]:} \textit{[Descripción del beneficio]}.
\end{enumerate}

\subsection{Beneficiarios Indirectos}

Adicionalmente, se beneficiarán de forma indirecta:

\begin{itemize}
    \item \textit{[Grupo indirecto 1]}: mediante \textit{[forma de beneficio]}
    \item \textit{[Grupo indirecto 2]}: a través de \textit{[mecanismo]}
\end{itemize}

\subsection{Alcance y Proyección Social}

El alcance social de esta investigación se proyecta a nivel:

\begin{itemize}
    \item \textbf{Local:} \textit{[Impacto en Chihuahua/comunidad específica]}
    \item \textbf{Regional:} \textit{[Posible extensión a otras regiones]}
    \item \textbf{Nacional:} \textit{[Aplicabilidad en México]}
    \item \textbf{Internacional:} \textit{[Potencial contribución global, si aplica]}
\end{itemize}

% ----------------------------------------------------------------------------
% CRITERIO 3: IMPLICACIONES PRÁCTICAS
% ----------------------------------------------------------------------------
\section{Implicaciones Prácticas}
\label{sec:implicaciones_practicas}

% Preguntas clave:
% - ¿Ayudará a resolver algún problema real?
% - ¿Tiene implicaciones trascendentales para problemas prácticos?

\subsection{Resolución de Problemas Reales}

Esta investigación ayudará a resolver el problema real de \textit{[problema específico]} que actualmente afecta a \textit{[población/contexto]}.

\textbf{Problemas prácticos que se abordan:}

\begin{enumerate}
    \item \textbf{Problema práctico 1:} \textit{[Descripción]}
    
    \textbf{Solución propuesta:} Los resultados permitirán \textit{[cómo ayudará a resolver]}.
    
    \item \textbf{Problema práctico 2:} \textit{[Descripción]}
    
    \textbf{Solución propuesta:} \textit{[Contribución específica]}.
\end{enumerate}

\subsection{Aplicaciones Prácticas Inmediatas}

Los resultados podrán aplicarse inmediatamente en:

\begin{itemize}
    \item \textit{[Aplicación práctica 1]}: En el contexto de \textit{[dónde/cuándo]}
    \item \textit{[Aplicación práctica 2]}: Para \textit{[qué propósito específico]}
    \item \textit{[Aplicación práctica 3]}: Beneficiando a \textit{[quiénes]}
\end{itemize}

\subsection{Trascendencia para Problemas Amplios}

Más allá del problema específico, esta investigación tiene implicaciones para:

\begin{itemize}
    \item \textit{[Problema más amplio 1]}
    \item \textit{[Problema más amplio 2]}
    \item \textit{[Problema más amplio 3]}
\end{itemize}

% ----------------------------------------------------------------------------
% CRITERIO 4: VALOR TEÓRICO
% ----------------------------------------------------------------------------
\section{Valor Teórico}
\label{sec:valor_teorico}

% Preguntas clave:
% - ¿Se llenará algún vacío de conocimiento?
% - ¿Se podrán generalizar resultados?
% - ¿Servirá para revisar/desarrollar/apoyar una teoría?
% - ¿Qué se sabrá que no se sabía antes?
% - ¿Se pueden sugerir hipótesis para futuros estudios?

\subsection{Vacíos de Conocimiento que se Llenarán}

Esta investigación llenará los siguientes vacíos de conocimiento:

\begin{enumerate}
    \item \textbf{Vacío 1:} Actualmente no se conoce \textit{[qué específicamente]}. Esta investigación aportará \textit{[qué nuevo conocimiento]}.
    
    \item \textbf{Vacío 2:} Existe poca evidencia sobre \textit{[aspecto]}. Los resultados proporcionarán \textit{[contribución]}.
\end{enumerate}

\subsection{Contribución al Desarrollo Teórico}

Los resultados de esta investigación contribuirán a:

\begin{itemize}
    \item \textbf{Revisar:} la teoría de \textit{[nombre]} al \textit{[cómo contribuirá]}
    \item \textbf{Desarrollar:} el concepto de \textit{[concepto]} mediante \textit{[aportación]}
    \item \textbf{Apoyar:} el modelo de \textit{[modelo]} al proporcionar \textit{[evidencia]}
\end{itemize}

\subsection{Conocimiento Nuevo que se Generará}

Con esta investigación se espera saber:

\begin{itemize}
    \item \textit{[Conocimiento nuevo 1]} que actualmente no se conoce
    \item \textit{[Conocimiento nuevo 2]} sobre el comportamiento de \textit{[variables]}
    \item \textit{[Conocimiento nuevo 3]} respecto a la relación entre \textit{[variables]}
\end{itemize}

\subsection{Generación de Hipótesis Futuras}

Los resultados permitirán sugerir:

\begin{itemize}
    \item Nuevas hipótesis sobre \textit{[tema]}
    \item Líneas de investigación futuras en \textit{[área]}
    \item Recomendaciones para estudios posteriores que aborden \textit{[aspecto]}
\end{itemize}

% ----------------------------------------------------------------------------
% CRITERIO 5: UTILIDAD METODOLÓGICA
% ----------------------------------------------------------------------------
\section{Utilidad Metodológica}
\label{sec:utilidad_metodologica}

% Preguntas clave:
% - ¿Ayudará a crear un nuevo instrumento?
% - ¿Contribuye a definir conceptos/variables/relaciones?
% - ¿Mejoras en la forma de experimentar?
% - ¿Sugiere cómo estudiar mejor una población?

\subsection{Aportaciones Metodológicas}

Esta investigación tiene utilidad metodológica porque:

\textbf{Instrumentos:}

\begin{itemize}
    \item Se desarrollará/validará \textit{[instrumento específico]} para \textit{[propósito]}
    \item Se adaptará \textit{[instrumento existente]} al contexto de \textit{[población/región]}
\end{itemize}

\textbf{Definiciones operacionales:}

\begin{itemize}
    \item Contribuirá a definir operacionalmente \textit{[concepto/variable]} en el contexto de \textit{[población]}
    \item Clarificará la relación entre \textit{[variables]}
\end{itemize}

\textbf{Mejoras metodológicas:}

\begin{itemize}
    \item Propondrá una mejora en la forma de medir \textit{[variable]} mediante \textit{[método]}
    \item Sugerirá una estrategia más eficiente para \textit{[proceso metodológico]}
\end{itemize}

\subsection{Replicabilidad y Transferibilidad}

La metodología empleada podrá ser:

\begin{itemize}
    \item \textbf{Replicada} en \textit{[otros contextos/poblaciones]}
    \item \textbf{Adaptada} para estudiar \textit{[problemas similares]}
    \item \textbf{Transferida} a investigaciones sobre \textit{[temas relacionados]}
\end{itemize}

% ----------------------------------------------------------------------------
% SÍNTESIS DE LA JUSTIFICACIÓN
% ----------------------------------------------------------------------------
\section{Síntesis de la Justificación}
\label{sec:sintesis_justificacion}

En síntesis, esta investigación se justifica plenamente porque:

\begin{enumerate}
    \item \textbf{Es conveniente} para \textit{[síntesis de conveniencia]}
    \item \textbf{Tiene relevancia social} al beneficiar a \textit{[síntesis de beneficiarios]}
    \item \textbf{Tiene implicaciones prácticas} al resolver \textit{[síntesis de problemas]}
    \item \textbf{Tiene valor teórico} al aportar \textit{[síntesis de conocimiento nuevo]}
    \item \textbf{Tiene utilidad metodológica} al desarrollar/validar \textit{[síntesis de aportación metodológica]}
\end{enumerate}

Por todo lo anterior, la realización de esta investigación es pertinente, necesaria y oportuna para \textit{[conclusión final]}.

% ============================================================================
% NOTAS FINALES:
% ============================================================================
%
% RECUERDA:
% 1. Justifica TU INVESTIGACIÓN, no el tema en general
% 2. Sé ESPECÍFICO en cada criterio
% 3. Responde TODAS las preguntas de cada criterio
% 4. Conecta con tu problema, pregunta y objetivos
% 5. Menciona beneficiarios CONCRETOS (números si es posible)
% 6. Las implicaciones prácticas deben ser REALES y ALCANZABLES
%
% ============================================================================
